\section*{Bab \ref{ch:g5:decimalops} --- Menghitung dengan Bilangan Desimal}

\begin{solution}{exo:g5:decimalops:1}
$45.7 + 8.93 = 54.63$; \quad
$16.25 + 3.75 = 20$; \quad
$0.86 + 12.4 = 13.26$.
\end{solution}

\begin{solution}{exo:g5:decimalops:2}
$9.4 - 2.72 = 6.68$ (periksa: $6.68 + 2.72 = 9.4$).

$20 - 13.45 = 6.55$ (periksa: $6.55 + 13.45 = 20$).

$6.03 - 5.9 = 0.13$.
\end{solution}

\begin{solution}{exo:g5:decimalops:3}
$6.2 \to 10$: $3.8$. \quad $3.75 \to 10$: $6.25$. \quad
$9.99 \to 10$: $0.01$.
\end{solution}

\begin{solution}{exo:g5:decimalops:4}
$72.4$; \quad $58$; \quad $3.45$; \quad $0.07$; \quad $300$.
\end{solution}

\begin{solution}{exo:g5:decimalops:5}
$4.05$ m $= 405$ cm; \quad $325$ cm $= 3.25$ m; \quad
$0.6$ kg $= 600$ g; \quad $85$ cL $= 0.85$ L.
\end{solution}

\begin{solution}{exo:g5:decimalops:6}
$6.2 \times 4$: taksiran $24$; tepatnya $24.8$.

$3.45 \times 8$: taksiran $28$ ($3.5 \times 8$); tepatnya $27.6$.

$12.5 \times 6$: taksiran $75$; tepatnya $75$.
\end{solution}

\begin{solution}{exo:g5:decimalops:7}
$7$ putaran: $0.4 \times 7 = 2.8$ km. $25$ putaran:
$0.4 \times 25 = 10$ km.
\end{solution}

\begin{solution}{exo:g5:decimalops:8}
Baguette: $1.15 \times 2 = 2.30$. Seluruhnya: $2.30 + 12.60 = 14.90$.
Kembalian: $20 - 14.90 = 5.10$.
\end{solution}

\begin{solution}{exo:g5:decimalops:9}
Satu paket: $1.5 \times 6 = 9$ L. Hari: $0.75 + 0.75 = 1.5$ L setiap
dua hari, jadi $9$ L bertahan $12$ hari (dua belas takaran $0.75$:
$12 \times 0.75 = 9$).
\end{solution}

\begin{solution}{exo:g5:decimalops:10}
Taksiran: $5.6$ lebih dari $5.5$, dan $5.5 \times 3 = 16.5$ --- jadi
\omterm{def:g2:mult:def}{hasil kalinya} pasti \emph{lebih} dari $16.5$, dan $15.18$ terlalu
kecil. Yang benar: $56 \times 3 = 168$, satu angka desimal, jadi
$5.6 \times 3 = 16.8$. Miko menempelkan kedua \omterm{def:g2:mult:def}{hasil kali} sebagiannya
($15$ dan $18$) berdampingan alih-alih menjumlahkannya sesuai nilai
tempatnya: $15 + 1.8 = 16.8$.
\end{solution}

\begin{solution}{exo:g5:decimalops:11}
\emph{1.} $65.4 + 64.9 + 65.0 + 66.1 = 261.4$ s.

\emph{2.} $261.4 > 260$: dia tidak memecahkan rekornya; dia
kelebihan $1.4$ s.

\end{solution}
